\begin{table}[ht]
\centering
\caption{Acurácia e ganho relativo do modelo LRM por perturbação ($T=0{,}6$). O símbolo (*) em \textbf{Ganho} indica variação estatisticamente significativa entre a condição Original e a perturbação (teste de McNemar pareado por exemplo, $p<0{,}05$; exato quando $b+c<25$).}
\label{tab:lrm_perturbacoes}
\sisetup{
  output-decimal-marker={,},
  table-format=-2.1,
  table-space-text-post={*}
}
\begin{tabular}{l
  S[table-format=2.1]
  S[table-format=-3.1, table-space-text-post={*}]
  S[table-format=3.0]
}
\toprule
\textbf{Perturbação} & {\textbf{Acurácia (\%)}} & {\textbf{Ganho (\%)}} & {\textbf{$n$}} \\
\midrule
Original & 57.0 & {--} & 521 \\
Linguagem Natural & 63.9 & +12.1* & 521 \\
Premissas Embaralhadas & 59.3 & +4.0 & 521 \\
Premissas Juntas (AND) & 58.2 & +2.0 & 521 \\
Ruído Irrelevante & 58.0 & +1.7 & 521 \\
Premissa Faltante & 55.8 & -2.0* & 197 \\
Duplicação de Premissas & 56.8 & -0.3 & 521 \\
Contradição Injetada & 46.8 & -17.8* & 521 \\
Negação Conclusão & 17.7 & -69.0* & 521 \\
\bottomrule
\end{tabular}
\end{table}